六足机器人模型预测控制（MPC）需要在高维混杂动力学上进行在线优化，同时保持关节物理解分支的正确性。我们的实验表明，带 warm start 的序列二次规划（SQP）流程可能收敛到任务空间上看似合理、但 tarsus 关节分支物理上错误的解；随后 runtime guard 虽可阻止危险指令下发，却会直接破坏前向推进。为解决这一问题，本文在最优控制问题中显式引入硬约束 \emph{TarsusBranchConstraint}，将左右 tarsus 的符号可行域直接写入求解器，并采用内点法（IPM）求解该受约束 MPC。与原有 SQP 路径不同，IPM 在求解过程中持续保持该不等式约束的活跃性，从而在可行域层面剔除错误分支，而非在优化后再进行修补。基于 Isaac Sim 闭环实验，在 18 执行器六足机器人上，本文方法在短程受控对照中将平均 branch-guard 触发次数从 275.9 次/轮降至 0.0 次/轮，将跟踪均方根误差（RMSE）从 0.242~m 降低到 0.175~m，并保持 4.82~ms 的平均 MPC 求解时间。在升级后的 endurance 配置下，5 次重复的 20~s 评测进一步获得了 0.306~m 的平均前向位移，且 guard 与 reject 均为零。进一步地，在 60~s endurance 评测中，本文方法实现了 0.638~m 的稳定前向位移，偏航漂移为 8.27$^\circ$，机身高度变化范围仅 0.011~m，平均求解时间保持在 4.63~ms。结果表明，求解器可见的物理分支硬约束是多足机器人可靠行走的关键，而 IPM 为实时、安全且具备长航时能力的六足 MPC 提供了一条可行路径。
